\section{Estudio de caso: Model-Based RL en manipulación diestra}

El ponente presentó un caso concreto en el que se usa model-based RL para que una mano robótica diestra manipule objetos.

\begin{figure}[H]
\centering
\includegraphics[width=0.9\textwidth]{slides-images/slide-35.jpg}
\caption{Resultados de los experimentos de ablación en simulación para la manipulación diestra\protect\footnotemark}
\end{figure}
\footnotetext{Slides, página 35.}

Hallazgos experimentales principales:
\begin{itemize}
    \item Se necesita una arquitectura de modelo suficientemente grande (2x250 o 2x500 unidades ocultas)
    \item Se requieren al menos 3 miembros del ensamble para obtener estimaciones de incertidumbre confiables
    \item El horizonte de planificación implica un compromiso: si es demasiado corto, no se ve lo bastante lejos en el futuro; si es demasiado largo, se acumula el error del modelo
    \item El controlador CEM mejorado (por ejemplo, Filtering + Reward Weighting) es mucho mejor que Random Shooting
\end{itemize}

\subsection{Cuándo usar Model-Based RL}

\begin{importantbox}{Escenarios de aplicación de Model-Based RL}
\begin{itemize}
    \item La recolección de datos es costosa (por ejemplo, con robots reales)
    \item El espacio de estados es de dimensión relativamente baja y predecible
    \item Se necesita adaptarse con rapidez a tareas nuevas (el modelo es transferible)
    \item La dinámica tiene cierta estructura que se puede aprender
\end{itemize}
Escenarios no adecuados: entornos con observaciones de alta dimensión (como imágenes) y dinámica altamente estocástica.
\end{importantbox}

\subsection{Resumen de la sección}

El caso de manipulación diestra muestra la ventaja de model-based RL en eficiencia de datos, así como la influencia decisiva de hiperparámetros como el ensamble y el horizonte de planificación en el desempeño.
